\section{Results}
\label{sec:results}

\begin{figure*}[t]
\vspace{-0.8em}
\centering
\includegraphics[width=0.72\textwidth]{sleep_main_results.png}
\caption{Sleep-event detection results.
A: event AP by matching tolerance.
B: matched GRU objective sweep comparing BDL kernels with segmentation losses.
C: matched backbone control comparing cross-entropy segmentation with BDL-Hard.
All panels use the same series-level folds, decoder search, and event-scoring protocol.}
\label{fig:sleep-main-results}
\vspace{-0.8em}
\end{figure*}


Figure~\ref{fig:sleep-main-results} shows the central empirical result: BDL improves the part of the precision--recall curve where event timing matters most.
At a one-minute tolerance, BDL-Hard reaches 0.219 AP, compared with 0.051 for cross-entropy.
At three minutes the comparison is 0.530 versus 0.283.
By thirty minutes the curves nearly meet, 0.838 versus 0.821, because both objectives often identify the sleep episode when the matching window is broad.
The separation at strict tolerances is therefore a localization effect, not merely better episode discovery.
This convergence at broad tolerances is also informative: BDL's advantage appears when the benchmark asks for a well-ranked detection rather than only the right episode.

One alternative explanation is class imbalance rather than BDL supervision itself.
With the GRU, split, scorer, and post-processing search fixed, BDL-Hard reaches 0.680 mAP, while cross-entropy, class-weighted cross-entropy, and focal loss reach 0.575, 0.544, and 0.525.
These segmentation controls change class balance without changing the interval-mask labels, so event scores are still recovered from state predictions after training.
BDL-Gaussian and BDL-Tolerance reach 0.660 and 0.605, showing that smoothing affects performance but does not account for the improvement by itself.
Among the BDL variants, the unsmoothed kernel is strongest, placing each annotated event at the recorded time.

The main architecture controls ask whether the pattern is tied to the recurrent encoder.
BDL-Hard improves over cross-entropy in each completed main architecture row, while the event scorer and validation-time decoder calibration are held fixed.
These comparisons replace the representation with recurrent and convolutional backbones designed for time-localized sequence prediction.
They also rule out a simpler explanation: segmentation is not weak only because the baseline recurrent encoder is mismatched to interval-mask prediction.
Across these controls the main change is the supervision and loss, while every trained model is still converted to ranked detections and scored by the same event-matching protocol.
Figure~\ref{fig:qualitative-window} illustrates the mechanism in a held-out sleep window: segmentation learns an interval mask and derives event scores afterward, while BDL directly scores the onset and wake-up events ranked by event AP.
Appendix Table~\ref{tab:paired-fold-consistency} shows the same direction on every matched sleep seed-fold comparison for the completed recurrent and convolutional controls.
\IfFileExists{results/generated/attention_unet_control.tex}{%
Appendix Table~\ref{tab:attention-unet-control} reports the attention-gated U-Net as a convolutional sensitivity check.
}{}
\IfFileExists{results/generated/transformer_strong_candidate_ready.tex}{%
A matched self-attention diagnostic is reported in Appendix Table~\ref{tab:transformer-strong-candidate}; in the current runs it is a negative control rather than evidence for the main architecture claim.
}{}

\IfFileExists{results/generated/seizure_highscore_ready.tex}{%
The high-capacity CHB-MIT comparison gives a second benchmark check under a
different signal modality and timing scale.
With a three-layer GRU and the standard one-to-sixty-second seizure scoring
protocol, BDL-Gaussian reaches 0.360 mAP at a two-second output stride,
compared with 0.252 for cross-entropy segmentation.
At a one-second output stride, BDL-Hard reaches 0.388 mAP compared with 0.287
for cross-entropy.
The preferred BDL kernel changes with output stride, which is consistent with
the resolution tradeoff in the stride heuristic: coarser outputs shorten the
sequence but make endpoint localization more dependent on smoothing and
decoder calibration.
The stricter endpoint analysis in the appendix separates boundary localization
from broader episode retrieval.
Across both stride settings, the main comparison is stable: after capacity and
output timeline are fixed, BDL remains above the matched segmentation baseline.
\IfFileExists{results/generated/seizure_highscore_strict3_ready.tex}{%
Under the stricter one-to-three-second endpoint analysis in the appendix,
BDL-Hard and BDL-Gaussian remain above cross-entropy, while the tolerance-shaped
kernel no longer helps.
}{}
}{}

Together, sleep and CHB-MIT support the same event-detection claim: when the benchmark ranks matched event times, aligning the training signal with event occurrences improves localization under matched scoring and post-processing settings.
The likelihood defines the training signal; kernel shape, output stride, representation, calibration, and decoder constraints still determine how well that signal becomes detections, especially on high-rate EEG.

The appendix uses the same logic to test alternative explanations.
If the gain came from a fortunate smoothing width, rare-class initialization, stronger encoder, or ordinary class-imbalance handling, then the kernel, prior, segmentation-reweighting, and architecture controls should close the gap under the same event scorer.
On sleep, these controls do not close the gap.
Other controls change the setting around the objective: CHB-MIT changes signal type and timing scale, point-event data remove paired endpoints, stride changes the output timeline, streaming and Transformer comparisons alter the available context and representation, and state reconstruction changes the quantity evaluated after training.
These studies define the scope of the method: BDL supplies event-scoring outputs for AP, while representation, output stride, context, calibration, and post-processing determine how accurately those scores become detections.

\begin{figure*}[t]
\vspace{-0.6em}
\centering
\includegraphics[width=0.60\textwidth]{sleep_prediction_example.png}
\caption{Held-out sleep window.
The shaded span is the annotated sleep interval; dashed green and blue lines mark onset and wake-up.
Segmentation learns a state sequence and derives boundary evidence from transitions, while BDL-Hard produces peaks in the score sequences at the scored event times.}
\label{fig:qualitative-window}
\vspace{-0.8em}
\end{figure*}

